% IRM：不同环境经同一表示后共享一个最优分类头。
\begin{tikzpicture}[
  env/.style={book node,fill=FigureInput!10,minimum width=23mm,minimum height=9mm},
  rep/.style={book node,fill=FigureModel!10,minimum width=25mm,minimum height=10mm},
  head/.style={book node,fill=FigureOutput!10,minimum width=24mm,minimum height=12mm},
  note/.style={font=\kaishu\scriptsize,text=BookMuted,align=center}
]
  \node[env] (e1) {环境$e=1$};
  \node[env,below=8mm of e1] (e2) {环境$e=2$};
  \node[env,below=8mm of e2] (e3) {环境$e=3$};
  \node[rep,right=15mm of e2] (phi) {共享表示$\Phi$};
  \node[head,right=18mm of phi] (w) {同一个头$w$};
  \node[book node,fill=FigureLoss!8,right=15mm of w] (risk) {各环境的期望风险$R^e$};

  \draw[book arrow] (e1)--(phi);
  \draw[book arrow] (e2)--(phi);
  \draw[book arrow] (e3)--(phi);
  \draw[book arrow] (phi)--(w);
  \draw[book arrow] (w)--(risk);

  \node[note,text=BookBlue,above=5mm of w] {原始约束：$w$在每个环境都最优};
  \node[note,text=FigureLoss,below=6mm of w] {IRMv1：在固定$w=1$处惩罚各环境梯度};
  \draw[FigureLoss,dashed,->] ([yshift=-4mm]risk.south) .. controls +(0,-8mm) and +(18mm,0) .. ([yshift=-4mm]w.south);
\end{tikzpicture}
